\scanpage{41}
\mspage{32}
% archive sheet number: 41
% inv: text | top | Exemple 2
\uline{Exemple 2}\quad Considérons l'application
% inv: display | upper third | varphi M to Figele L
\[
\varphi:M\longrightarrow\unsure{\mathrm{Fig\acute{e}l}|\mathrm{Fig\acute{e}le}}(L)\qquad X\longrightarrow\unsure{\mathrm{mulomb}|\mathrm{muloms}}(\unsure{X})
\]
% inv: deletion | upper third | On sait que c'est
\deleted{On sait que c'est}
% inv: text | upper third | Ce n'est défini
Ce n'est défini \inserted{(comme appli. à val. \unsure{ds|dans} \uuline{Figél})} que si on suppose qu'on a Mag L 1, i.e. $\forall X\in M$, $\mathrm{omb}(X)^\circ\ne\varnothing$.

% inv: text | upper middle | Dans ce cas
Dans ce cas, on a \deleted{bien}
% inv: display | upper middle | cancelled implication
\[
\deleted{X\mathrel{|\circ|}Y\Longrightarrow\mathrm{omb}(X)}
\]
% inv: display | middle | mulomb supports
\[
\begin{aligned}
|\mathrm{mulomb}(X)|&=\mathrm{omb}(X)\\
|\mathrm{mulomb}(X)|^\circ&=\mathrm{omb}(X)^\circ
\end{aligned}
\]
% inv: text | middle | et
et
% inv: display | middle | disjoint mulomb equivalence
\[
\mathrm{mulomb}(X)\mathrel{|\circ|}\mathrm{mulomb}(Y)\Longleftrightarrow\mathrm{omb}(X)^\circ\cap\mathrm{omb}(Y)^\circ=\unsure{\varnothing}
\]
% inv: text | middle lower | Donc on a bien
Donc on a bien
% inv: display | middle lower | implication to mulomb Y
\[
X\mathrel{|\circ|}Y\Longrightarrow\mathrm{mulomb}(X)\mathrel{|\circ|}\mathrm{mulomb}\,Y.
\]

% inv: deletion | lower middle | Cela nous montre
\deleted{Cela nous \unsure{montre}}
% inv: text | lower third | Ces deux exemples
Ces deux exemples nous montrent en tout cas que la notion de comorphisme proposée à la page 27, n'est pas trop idiote. (même s'il convient sans doute de la généraliser). Troisième exemple :

% inv: text | bottom third | Exemple 3
\uline{Exemple 3}\quad $f:\unsure{E|\mathbb{E}}'\longrightarrow E$ application ensembliste surjective. Elle induit un comorphisme injectif des magasins ensemblistes associés
% inv: display | bottom | Figel E to Figel E prime
\[
\mathrm{Figel}(E)\longrightarrow\mathrm{Figel}(E')
\]
